%&these

\makeindex

\begin{document}

\title{
    {Approche Géométrique pour la Conception de Codes à Effacement Appliqués
    dans le Stockage Distribué}\\
    {\large IRCCyN}\\
    {\includegraphics[width=0.5\textwidth]{img/logo_irccyn.pdf}}
}
\author{Dimitri Pertin}
\date {}

\maketitle

\newpage

\tableofcontents

\part{Code à effacement haute performance}

\chapter{Principe et évaluation des codes à effacement}

\input{chapters/chapter01.tex}

\chapter{Conception et relation de codes définis par géométrie discrète}

La Mojette\index{mojette} ça poutre.

\input{chapters/chapter02.tex}

\part{Application au stockage distribué}

\chapter{Redondance des systèmes de stockage distribués par codage à
effacement}

\input{chapters/chapter04.tex}

\chapter{Application au système de fichiers distribués RozoFS}

\input{chapters/chapter05.tex}

\chapter{Calcul distribué de nouveaux mots encodés}

\input{chapters/chapter03.tex}

\newpage

\printbibliography[
%    segment=1
]

\printindex

\end{document}
